% GENERATED FILE. Edit canonical lessons, then run npm run generate:books.
% canonical-source-hash: fnv1a64:c5be3c72a0c0c622
% canonical-lessons: MW-C255-tabiyat, MW-C255-bigarno, MW-C255-sudharno-better, MW-C255-asu, MW-C255-kapno

\chapter{Better or Worse}
\label{ch:mw-better-or-worse}
\hlchaptermodality{255}

\begin{chapteropening}
\textbf{By the end of this chapter:} \emph{I can say whether my health got better.}

Complete the last lesson of chapter 255: I can say whether my health got better.
\end{chapteropening}

\section[\texorpdfstring{\mw{तबियत} (tabiyat)}{तबियत (tabiyat)}]{\mw{तबियत} (tabiyat) — one's health; how one feels}

\label{lesson:MW-C255-tabiyat}



\begin{quote}
Before the new one: say the Marwadi for rest, then the Marwadi for sluggishness.
\end{quote}

\subsection*{You'll want to know: \mw{तबियत}}
\textbf{\mw{तबियत}} — \emph{tabiyat} — \textquotedblleft{}one's health; how one feels\textquotedblright{}.

\begin{quote}
\mw{आज} \mw{म्हारी} \mw{तबियत} \mw{ठीक} \mw{है।}
\end{quote}

\begin{quote}
\emph{āj mhārī tabiyat ṭhīk hai.}
\end{quote}

\begin{quote}
I am feeling well today.
\end{quote}

A feminine noun.

\begin{cousinweb}
\textbf{\mw{तबियत}} (\emph{tabiyat}) is usually traced to Persian \emph{tabīat}, \textquotedblleft{}nature, temperament\textquotedblright{}; in Marwari it is how one feels in body.
\end{cousinweb}

\begin{grammarlens}[title={What you've built}]
One more word for this chapter.
\end{grammarlens}

\subsection*{Your turn}
\begin{itemize}
\raggedright
  \item \emph{Say it:} \emph{tabiyat}
  \item \emph{Say it:} \emph{tabiyat}, once more
  \item \emph{Say it:} \emph{āj mhārī tabiyat ṭhīk hai.} — I am feeling well today.
  \item \emph{Recall:} say the Marwadi for rest, then the Marwadi for sluggishness, then say \emph{tabiyat} again
\end{itemize}

\begin{tcolorbox}[breakable,colback=teal!4,colframe=teal!35!black,title={Before you move on}]
What does \mw{तबियत} mean? (\textquotedblleft{}One's health; how one feels\textquotedblright{}.) Say it once more.
\end{tcolorbox}
\section[\texorpdfstring{\mw{बिगड़णो} (bigaṛṇo)}{बिगड़णो (bigaṛṇo)}]{\mw{बिगड़णो} (bigaṛṇo) — to get worse; to go wrong}

\label{lesson:MW-C255-bigarno}



\begin{quote}
Before the new one: say the Marwadi for sluggishness, then the Marwadi for one's health.
\end{quote}

\subsection*{You'll want to know: \mw{बिगड़णो}}
\textbf{\mw{बिगड़णो}} — \emph{bigaṛṇo} — \textquotedblleft{}to get worse; to go wrong\textquotedblright{}.

\begin{quote}
\mw{काल} \mw{म्हारी} \mw{तबियत} \mw{बिगड़ी।}
\end{quote}

\begin{quote}
\emph{kāl mhārī tabiyat bigaṛī.}
\end{quote}

\begin{quote}
Yesterday I felt worse.
\end{quote}

Past (it …): \mw{वो} \mw{दूध} \mw{बिगड़्यो}, \emph{vo dūdh bigaṛyo}.

\begin{cousinweb}
\textbf{\mw{बिगड़णो}} (\emph{bigaṛṇo}) is inherited from Sanskrit \emph{vighaṭate}, \textquotedblleft{}is marred\textquotedblright{}. CDIAL gives no Marwari form in this entry; it lists Hindi \emph{bigaṛnā}, \textquotedblleft{}to be spoilt\textquotedblright{} (CDIAL 11673).
\end{cousinweb}

\begin{grammarlens}[title={What you've built}]
One more word for this chapter.
\end{grammarlens}

\subsection*{Your turn}
\begin{itemize}
\raggedright
  \item \emph{Say it:} \emph{bigaṛṇo}
  \item \emph{Say it:} \emph{bigaṛṇo}, once more
  \item \emph{Say it:} \emph{kāl mhārī tabiyat bigaṛī.} — Yesterday I felt worse.
  \item \emph{Recall:} say the Marwadi for sluggishness, then the Marwadi for one's health, then say \emph{bigaṛṇo} again
\end{itemize}

\begin{tcolorbox}[breakable,colback=teal!4,colframe=teal!35!black,title={Before you move on}]
What does \mw{बिगड़णो} mean? (\textquotedblleft{}To get worse; to go wrong\textquotedblright{}.) Say it once more.
\end{tcolorbox}
\section[\texorpdfstring{\mw{सुधरणो} (sudharṇo)}{सुधरणो (sudharṇo)}]{\mw{सुधरणो} (sudharṇo) — to get better; to improve}

\label{lesson:MW-C255-sudharno-better}



\begin{quote}
Before the new one: say the Marwadi for one's health, then the Marwadi for to get worse.
\end{quote}

\subsection*{You'll want to know: \mw{सुधरणो}}
\textbf{\mw{सुधरणो}} — \emph{sudharṇo} — \textquotedblleft{}to get better; to improve\textquotedblright{}.

\begin{quote}
\mw{आज} \mw{तबियत} \mw{सुधरी।}
\end{quote}

\begin{quote}
\emph{āj tabiyat sudharī.}
\end{quote}

\begin{quote}
Today I am feeling better.
\end{quote}

Past (it …): \mw{वो} \mw{मौसम} \mw{सुधर्यो}, \emph{vo mausam sudharyo}.

\begin{cousinweb}
\textbf{\mw{सुधरणो}} (\emph{sudharṇo}) is inherited from a reconstructed \emph{*śuddhakāra}, \textquotedblleft{}making clean\textquotedblright{}. CDIAL gives no Marwari form in this entry; it lists Hindi \emph{sudharnā}, \textquotedblleft{}to be mended\textquotedblright{} (CDIAL 12521).
\end{cousinweb}

\begin{grammarlens}[title={What you've built}]
One more word for this chapter.
\end{grammarlens}

\subsection*{Your turn}
\begin{itemize}
\raggedright
  \item \emph{Say it:} \emph{sudharṇo}
  \item \emph{Say it:} \emph{sudharṇo}, once more
  \item \emph{Say it:} \emph{āj tabiyat sudharī.} — Today I am feeling better.
  \item \emph{Recall:} say the Marwadi for one's health, then the Marwadi for to get worse, then say \emph{sudharṇo} again
\end{itemize}

\begin{tcolorbox}[breakable,colback=teal!4,colframe=teal!35!black,title={Before you move on}]
What does \mw{सुधरणो} mean? (\textquotedblleft{}To get better; to improve\textquotedblright{}.) Say it once more.
\end{tcolorbox}
\section[\texorpdfstring{\mw{आंसू} (ā̃sū)}{आंसू (ā̃sū)}]{\mw{आंसू} (ā̃sū) — a tear (from the eye)}

\label{lesson:MW-C255-asu}



\begin{quote}
Before the new one: say the Marwadi for to get worse, then the Marwadi for to get better.
\end{quote}

\subsection*{You'll want to know: \mw{आंसू}}
\textbf{\mw{आंसू}} — \emph{ā̃sū} — \textquotedblleft{}a tear (from the eye)\textquotedblright{}.

\begin{quote}
\mw{छोरी} \mw{रोवै} \mw{है}, \mw{आंख} \mw{मांय} \mw{आंसू} \mw{है।}
\end{quote}

\begin{quote}
\emph{chhorī rovai hai, ā̃kh mā̃y ā̃sū hai.}
\end{quote}

\begin{quote}
The girl is crying; there are tears in her eyes.
\end{quote}

A masculine noun.

\begin{cousinweb}
\textbf{\mw{आंसू}} (\emph{ā̃sū}) is inherited from Sanskrit \emph{aśru}, \textquotedblleft{}a tear\textquotedblright{}. CDIAL gives no Marwari form in this entry; it lists Hindi \emph{ā̃sū} (CDIAL 919).
\end{cousinweb}

\begin{grammarlens}[title={What you've built}]
One more word for this chapter.
\end{grammarlens}

\subsection*{Your turn}
\begin{itemize}
\raggedright
  \item \emph{Say it:} \emph{ā̃sū}
  \item \emph{Say it:} \emph{ā̃sū}, once more
  \item \emph{Say it:} \emph{chhorī rovai hai, ā̃kh mā̃y ā̃sū hai.} — The girl is crying; there are tears in her eyes.
  \item \emph{Recall:} say the Marwadi for to get worse, then the Marwadi for to get better, then say \emph{ā̃sū} again
\end{itemize}

\begin{tcolorbox}[breakable,colback=teal!4,colframe=teal!35!black,title={Before you move on}]
What does \mw{आंसू} mean? (\textquotedblleft{}A tear (from the eye)\textquotedblright{}.) Say it once more.
\end{tcolorbox}
\section[\texorpdfstring{\mw{कांपणो} (kā̃pṇo)}{कांपणो (kā̃pṇo)}]{\mw{कांपणो} (kā̃pṇo) — to tremble; to shiver}

\label{lesson:MW-C255-kapno}



\begin{quote}
Before the new one: say the Marwadi for to get better, then the Marwadi for a tear (from the eye).
\end{quote}

\subsection*{You'll want to know: \mw{कांपणो}}
\textbf{\mw{कांपणो}} — \emph{kā̃pṇo} — \textquotedblleft{}to tremble; to shiver\textquotedblright{}.

\begin{quote}
\mw{ठंडी} \mw{मांय} \mw{हूं} \mw{कांपूं} \mw{हूं।}
\end{quote}

\begin{quote}
\emph{ṭhaṇḍī mā̃y hū̃ kā̃pū̃ hū̃.}
\end{quote}

\begin{quote}
I shiver in the cold.
\end{quote}

Past (I …): \mw{हूं} \mw{रात} \mw{मांय} \mw{कांप्यो}, \emph{hū̃ rāt mā̃y kā̃pyo}.

\begin{cousinweb}
\textbf{\mw{कांपणो}} (\emph{kā̃pṇo}) is inherited from Sanskrit \emph{kampate}, \textquotedblleft{}trembles\textquotedblright{}. CDIAL gives no Marwari form in this entry; it lists Hindi \emph{kā̃pnā} (CDIAL 2767).
\end{cousinweb}

\begin{grammarlens}[title={What you've built}]
One more word for this chapter.
\end{grammarlens}

\subsection*{Your turn}
\begin{itemize}
\raggedright
  \item \emph{Say it:} \emph{kā̃pṇo}
  \item \emph{Say it:} \emph{kā̃pṇo}, once more
  \item \emph{Say it:} \emph{ṭhaṇḍī mā̃y hū̃ kā̃pū̃ hū̃.} — I shiver in the cold.
  \item \emph{Recall:} say the Marwadi for to get better, then the Marwadi for a tear (from the eye), then say \emph{kā̃pṇo} again
\end{itemize}

\begin{tcolorbox}[breakable,colback=teal!4,colframe=teal!35!black,title={Before you move on}]
What does \mw{कांपणो} mean? (\textquotedblleft{}To tremble; to shiver\textquotedblright{}.) Say it once more.
\end{tcolorbox}
